\chapter{Conditional Probability, Independence and Bayes' Rule}
\label{ch:conditional}

Probabilities change when we learn something. The probability that a card is
the ace of spades is $1/52$; once we are told that it is black, it is $1/26$.
This chapter makes that idea precise, and develops the three tools that follow
from it: the multiplication rule, the law of total probability, and Bayes' rule.
We then define independence, the situation where learning one fact tells us
nothing about another.

\section{Conditional probability}

Suppose $\Omega$ is finite with equally likely outcomes, and we are told that
the event $B$ has occurred. The outcome is now equally likely to be any element
of $B$, and the event $A$ occurs exactly when the outcome lies in $A\cap B$. So
the natural updated probability is
\[
\frac{\abs{A\cap B}}{\abs{B}}=\frac{\abs{A\cap B}/\abs\Omega}{\abs B/\abs\Omega}
=\frac{\Prob(A\cap B)}{\Prob(B)}.
\]
The right-hand side makes sense on any probability space, which motivates the
following definition.

\begin{definition}[Conditional probability]
Let $B$ be an event with $\Prob(B)>0$. The \emph{conditional probability of $A$
given $B$} is
\[
\Prob(A\given B)=\frac{\Prob(A\cap B)}{\Prob(B)}.
\]
\end{definition}

\begin{theorem}[Conditioning gives a probability measure]\label{thm:condmeasure}
If $\Prob(B)>0$, the function $A\mapsto\Prob(A\given B)$ is a probability
measure on $(\Omega,\F)$. Consequently every result of Chapter~\ref{ch:axioms}
holds with $\Prob(\cdot)$ replaced by $\Prob(\cdot\given B)$.
\end{theorem}
\begin{proof}
Non-negativity is clear. Normalisation: $\Prob(\Omega\given B)=\Prob(B)/\Prob(B)=1$.
For countable additivity, let $A_1,A_2,\dots$ be pairwise disjoint. The events
$A_i\cap B$ are also pairwise disjoint, with
$\bigl(\bigcup_iA_i\bigr)\cap B=\bigcup_i(A_i\cap B)$, so
\[
\Prob\Bigl(\bigcup_iA_i\given B\Bigr)
=\frac{\sum_i\Prob(A_i\cap B)}{\Prob(B)}=\sum_i\Prob(A_i\given B).\qedhere
\]
\end{proof}

\begin{nuance}{conditioning on events of probability zero}
The definition requires $\Prob(B)>0$. Conditioning on an event of probability
zero is not simply ``undefined'': for continuous random variables one
conditions on $\{X=x\}$, which has probability zero, and obtains perfectly
meaningful answers. This requires a different construction (conditional
densities in Chapter~\ref{ch:jointcont} and conditional expectation in
Chapter~\ref{ch:condexp}, made fully rigorous in Chapter~\ref{ch:measure}). Until then, all conditioning events have positive
probability.
\end{nuance}

\begin{pitfall}{$\Prob(A\given B)$ and $\Prob(B\given A)$ are different}
In general $\Prob(A\given B)\ne\Prob(B\given A)$. The two quantities share the
numerator $\Prob(A\cap B)$ but have different denominators. For example, almost
every Olympic sprinter is a good runner, but few good runners are Olympic
sprinters. Confusing the two is the single most common error in applied
probability; the examples at the end of Section~\ref{sec:bayes} show its
consequences.
\end{pitfall}

\section{The multiplication rule and total probability}

Rearranging the definition gives a rule for the probability of a joint event.

\begin{theorem}[Multiplication (chain) rule]\label{thm:chain}
If $\Prob(A_1\cap\dots\cap A_{n-1})>0$, then
\[
\Prob(A_1\cap\dots\cap A_n)
=\Prob(A_1)\,\Prob(A_2\given A_1)\,\Prob(A_3\given A_1\cap A_2)\cdots
\Prob(A_n\given A_1\cap\dots\cap A_{n-1}).
\]
\end{theorem}
\begin{proof}
Write each conditional probability as a ratio. The product telescopes:
\[
\Prob(A_1)\cdot\frac{\Prob(A_1\cap A_2)}{\Prob(A_1)}\cdot
\frac{\Prob(A_1\cap A_2\cap A_3)}{\Prob(A_1\cap A_2)}\cdots
\frac{\Prob(A_1\cap\dots\cap A_n)}{\Prob(A_1\cap\dots\cap A_{n-1})}
=\Prob(A_1\cap\dots\cap A_n).
\]
All denominators are positive: for each $k\le n-1$ we have
$A_1\cap\dots\cap A_k\supseteq A_1\cap\dots\cap A_{n-1}$, so by monotonicity
(Proposition~\ref{prop:basic}(iv)) its probability is at least that of
$A_1\cap\dots\cap A_{n-1}$, which is positive by hypothesis.
\end{proof}

\begin{example}[Drawing without replacement]
An urn has $5$ red and $3$ blue balls. Two are drawn without replacement. Let
$R_i$ be ``the $i$th ball is red''. Then
\[
\Prob(R_1\cap R_2)=\Prob(R_1)\Prob(R_2\given R_1)=\frac58\cdot\frac47=\frac5{14}.
\]
\end{example}

\begin{theorem}[Law of total probability]\label{thm:lotp}
Let $B_1,B_2,\dots$ be a finite or countable partition of $\Omega$ with
$\Prob(B_i)>0$ for every $i$. Then for every event $A$,
\[
\Prob(A)=\sum_i\Prob(A\given B_i)\,\Prob(B_i).
\]
\end{theorem}
\begin{proof}
The events $A\cap B_i$ are pairwise disjoint and their union is $A$ (because the
$B_i$ cover $\Omega$). By countable additivity,
$\Prob(A)=\sum_i\Prob(A\cap B_i)$, and by the definition of conditional
probability $\Prob(A\cap B_i)=\Prob(A\given B_i)\Prob(B_i)$.
\end{proof}

The law of total probability embodies a strategy: split a hard problem into
cases according to a partition, solve each case, and average with weights
$\Prob(B_i)$. It is the starting point for \emph{first-step analysis}, which we
use extensively for random walks and Markov chains.

\begin{example}[Two urns]
Urn~I has $2$ red and $3$ blue balls; urn~II has $4$ red and $1$ blue. An urn is
chosen with equal probability, then a ball is drawn at random from it. With
$B_1,B_2$ the events ``urn~I'', ``urn~II'',
\[
\Prob(\text{red})=\frac25\cdot\frac12+\frac45\cdot\frac12=\frac35 .
\]
\end{example}

\section{Bayes' rule}\label{sec:bayes}

\begin{theorem}[Bayes' rule]\label{thm:bayes}
Let $B_1,B_2,\dots$ be a partition of $\Omega$ with $\Prob(B_i)>0$, and let $A$
be an event with $\Prob(A)>0$. Then for each $j$,
\[
\Prob(B_j\given A)=\frac{\Prob(A\given B_j)\Prob(B_j)}{\sum_i\Prob(A\given B_i)\Prob(B_i)} .
\]
\end{theorem}
\begin{proof}
By definition, $\Prob(B_j\given A)=\Prob(A\cap B_j)/\Prob(A)$ and
$\Prob(A\cap B_j)=\Prob(A\given B_j)\Prob(B_j)$. Replace $\Prob(A)$ in the
denominator using Theorem~\ref{thm:lotp}.
\end{proof}

In words: the \emph{posterior} probability of a hypothesis $B_j$ after seeing the
evidence $A$ is proportional to the \emph{prior} $\Prob(B_j)$ times the
\emph{likelihood} $\Prob(A\given B_j)$, and the denominator simply normalises.
For two hypotheses it is convenient to use odds.

\begin{corollary}[Odds form]
For events $B$ and $A$ with $0<\Prob(B)<1$ and $\Prob(A)>0$,
\[
\frac{\Prob(B\given A)}{\Prob(B^c\given A)}
=\frac{\Prob(B)}{\Prob(B^c)}\cdot\frac{\Prob(A\given B)}{\Prob(A\given B^c)} .
\]
That is, posterior odds $=$ prior odds $\times$ likelihood ratio.
\end{corollary}
\begin{proof}
Apply the ratio form of Bayes' rule to $B$ and $B^c$; the common denominator
$\Prob(A)$ cancels.
\end{proof}

\begin{example}[A medical test and the base rate]\label{ex:medtest}
A disease affects $1\%$ of a population. A test returns positive for $99\%$ of
people with the disease (sensitivity) and for $5\%$ of people without it (false
positive rate). A random person tests positive. What is the probability they
have the disease? Let $D$ be ``has the disease'' and $+$ ``tests positive''.
Then
\[
\Prob(D\given +)=\frac{0.99\cdot0.01}{0.99\cdot0.01+0.05\cdot0.99}
=\frac{0.0099}{0.0099+0.0495}=\frac{0.0099}{0.0594}=\frac16\approx0.167 .
\]
Even though the test is accurate, only about one positive in six is a true case,
because healthy people vastly outnumber sick ones. In odds form: the prior odds
are $1:99$ and the likelihood ratio is $0.99/0.05=19.8$, so the posterior odds
are $19.8:99=1:5$, which is $1/6$ again.
\end{example}

\begin{pitfall}{ignoring the base rate}
Many people answer ``about $95$--$99\%$'' in Example~\ref{ex:medtest}, treating
$\Prob(+\given D)$ as if it were $\Prob(D\given +)$. The correct answer needs
the prior $\Prob(D)$. A useful sanity check: imagine $10{,}000$ people.
About $100$ have the disease and $99$ of them test positive; of the $9{,}900$
healthy, about $495$ test positive. Among about $594$ positives, $99$ are
sick.
\end{pitfall}

\begin{example}[The prosecutor's fallacy]
Suppose a DNA profile matches a randomly chosen innocent person with probability
$10^{-6}$, and a crime is committed in a city of $10^7$ people, one of whom is
the culprit and matches with certainty. A suspect is found only because the
profile matched. The probability of the evidence if the suspect is innocent is
$10^{-6}$, but that is not the probability of innocence. Among the roughly
$10^7\cdot10^{-6}=10$ innocent people who would also match, plus the guilty
one, only about $1$ in $11$ matches is the culprit, if no other evidence
distinguishes them. Treating $\Prob(\text{evidence}\given\text{innocent})$ as
$\Prob(\text{innocent}\given\text{evidence})$ is exactly the confusion in the
pitfall above.
\end{example}

\begin{example}[The Monty Hall problem]\label{ex:monty}
Behind one of three doors is a car; behind the others, goats. You pick door~1.
The host, who knows where the car is, opens one of the other two doors,
\emph{always revealing a goat}, and if both remaining doors hide goats he picks
one at random. Should you switch?

Let $C_i$ be ``car behind door $i$'', each with probability $1/3$, and let $H$ be
the event that the host opens door~3. Then
$\Prob(H\given C_1)=1/2$, $\Prob(H\given C_2)=1$ (he must avoid the car), and
$\Prob(H\given C_3)=0$. By total probability $\Prob(H)=\frac13(\frac12+1+0)
=\frac12$, and by Bayes' rule
\[
\Prob(C_1\given H)=\frac{\frac12\cdot\frac13}{\frac12}=\frac13,\qquad
\Prob(C_2\given H)=\frac{1\cdot\frac13}{\frac12}=\frac23 .
\]
Switching to door~2 wins with probability $2/3$.
\end{example}

\begin{nuance}{the answer depends on the host's protocol}
Monty Hall is famous because its answer depends on a sentence that is easy to
skip: the host \emph{knows} where the car is and \emph{always} reveals a goat.
Suppose instead that the host opens one of the two unchosen doors at random, and
we observe that he happened to reveal a goat. Then
$\Prob(\text{goat shown}\given C_1)=1$ but
$\Prob(\text{goat shown}\given C_2)=\Prob(\text{goat shown}\given C_3)=1/2$,
so the probability that your original door wins, given that a goat was shown, is
\[
\frac{\frac13\cdot1}{\frac13\cdot1+\frac23\cdot\frac12}=\frac12,
\]
and switching gives no advantage. The same observed event (``door 3 shows a
goat'') has different probabilities under different protocols. In general, to
condition correctly you must model \emph{how the information arrived}, not just
the information.
\end{nuance}

\begin{nuance}{the two-children problem}
``A family has two children. At least one is a boy. What is the probability
both are boys?'' With $\Omega=\{BB,BG,GB,GG\}$ equally likely, conditioning on
``at least one boy'' leaves $\{BB,BG,GB\}$, so the answer is $1/3$. If instead we
are told that the \emph{older} child is a boy, only $\{BB,BG\}$ remain and the
answer is $1/2$. If we are told only that we met one of the children and it was
a boy, the answer depends on how that child was chosen. As with Monty Hall, the
wording specifies the information-gathering process, and different processes
give different answers.
\end{nuance}

\section{Independence}

Informally, $A$ and $B$ are independent if learning that $B$ occurred does not
change the probability of $A$: $\Prob(A\given B)=\Prob(A)$. Writing this as
$\Prob(A\cap B)=\Prob(A)\Prob(B)$ removes the need for $\Prob(B)>0$ and makes the
definition symmetric.

\begin{definition}[Independence of two events]
Events $A$ and $B$ are \emph{independent} if $\Prob(A\cap B)=\Prob(A)\Prob(B)$.
\end{definition}

\begin{proposition}\label{prop:indcomp}
If $A$ and $B$ are independent, then so are $A$ and $B^c$, $A^c$ and $B$, and
$A^c$ and $B^c$.
\end{proposition}
\begin{proof}
$\Prob(A\cap B^c)=\Prob(A)-\Prob(A\cap B)=\Prob(A)-\Prob(A)\Prob(B)
=\Prob(A)\Prob(B^c)$. The other two follow by applying this result again.
\end{proof}

\begin{pitfall}{independent is not the same as disjoint}
If $A$ and $B$ are disjoint and both have positive probability, they are
\emph{dependent}: $\Prob(A\cap B)=0\ne\Prob(A)\Prob(B)$. Indeed, learning that
$B$ occurred tells you that $A$ did not. Independence is about information, not
about set overlap.
\end{pitfall}

For more than two events, the right definition requires every sub-collection to
factor.

\begin{definition}[Mutual independence]
Events $A_1,\dots,A_n$ are \emph{mutually independent} if for every
sub-collection of indices $i_1<\dots<i_k$, $k\ge2$,
\[
\Prob(A_{i_1}\cap\dots\cap A_{i_k})=\Prob(A_{i_1})\cdots\Prob(A_{i_k}).
\]
This is $2^n-n-1$ conditions. An infinite family is mutually independent if
every finite sub-family is.
\end{definition}

\begin{example}[Pairwise independent but not mutually independent]\label{ex:pairwise}
Toss two fair coins. Let $A$ = ``first is heads'', $B$ = ``second is heads'',
$C$ = ``both tosses agree''. Then $\Prob(A)=\Prob(B)=\Prob(C)=1/2$. Each pairwise
intersection is $\{HH\}$ with probability $1/4=\frac12\cdot\frac12$, so the events
are pairwise independent. But $A\cap B\cap C=\{HH\}$ has probability $1/4$, not
$1/8$. Knowing any two of the events determines the third.
\end{example}

\begin{nuance}{checking only pairs is not enough}
Pairwise independence is strictly weaker than mutual independence. Many results
(laws of large numbers, concentration inequalities) hold under pairwise
independence, but others (such as the product formula for the probability that
none of the events occurs) require mutual independence. Always state which one
you are using.
\end{nuance}

\begin{example}[Independent trials]\label{ex:trials}
Independent trials each succeeding with probability $p$ model repeated
experiments such as coin tosses. The probability of exactly $k$ successes in $n$
trials is obtained by noting that each particular sequence with $k$ successes
has probability $p^k(1-p)^{n-k}$ (by mutual independence), and there are
$\binom nk$ such sequences (Theorem~\ref{thm:choose}):
\[
\Prob(k\text{ successes})=\binom nk p^k(1-p)^{n-k}.
\]
This is the \emph{binomial distribution}, studied in Chapter~\ref{ch:dist}.
\end{example}

\section{Conditional independence}

\begin{definition}
Given an event $C$ with $\Prob(C)>0$, events $A$ and $B$ are \emph{conditionally
independent given $C$} if $\Prob(A\cap B\given C)=\Prob(A\given C)\Prob(B\given C)$.
\end{definition}

This is independence under the probability measure $\Prob(\cdot\given C)$,
which Theorem~\ref{thm:condmeasure} shows is a legitimate measure.

\begin{example}[Conditionally independent but dependent]\label{ex:condind}
Two coins: coin~1 shows heads with probability $0.9$, coin~2 with probability
$0.1$. Choose a coin with equal probability, then toss it twice. Given the coin,
the tosses are independent. Unconditionally, let $H_1,H_2$ be ``heads on toss
1/2''. Then $\Prob(H_1)=\Prob(H_2)=\frac12(0.9+0.1)=0.5$ but
\[
\Prob(H_1\cap H_2)=\tfrac12(0.9^2)+\tfrac12(0.1^2)=0.41\ne0.25 .
\]
The first toss is informative about which coin we hold, and therefore about the
second toss.
\end{example}

\begin{nuance}{independence and conditional independence are logically unrelated}
Conditional independence given $C$ does not imply independence (previous
example), and independence does not imply conditional independence. For
instance, two independent fair coin tosses become dependent when conditioned
on the event ``at least one head''. Learning that a third variable is
related to both can create or destroy dependence. This phenomenon underlies
Simpson's paradox below, and appears in Bayesian networks and causal inference.
\end{nuance}

\begin{example}[Simpson's paradox]\label{ex:simpson}
A well-known data set (Charig et al., 1986) compares two kidney-stone
treatments by success counts:
\begin{center}
\begin{tabular}{@{}lccc@{}}
\toprule
 & Small stones & Large stones & Overall\\
\midrule
Treatment A & $81/87\ (93\%)$ & $192/263\ (73\%)$ & $273/350\ (78\%)$\\
Treatment B & $234/270\ (87\%)$ & $55/80\ (69\%)$ & $289/350\ (83\%)$\\
\bottomrule
\end{tabular}
\end{center}
Treatment A is better within each stone size, yet B has the higher overall rate.
Law of total probability explains it: the overall success rate is
$\Prob(S\given T)=\sum_z\Prob(S\given T,Z=z)\Prob(Z=z\given T)$, a weighted
average of the group rates, and the weights $\Prob(Z=z\given T)$ differ: A was
used mostly on large (harder) stones, B mostly on small ones. No arithmetic is
wrong; the aggregate comparison simply mixes the effect of the treatment with
the effect of case difficulty.
\end{example}

\section*{Exercises}
\addcontentsline{toc}{section}{Exercises}

\begin{enumerate}
\item Two cards are dealt from a standard deck. Find the probability that the
second card is an ace given that the first is an ace, and then the probability
that the second is an ace given that the first is not an ace. Use total
probability to show that the unconditional probability that the second card is
an ace is $1/13$.
\item Prove that if $\Prob(A\given B)>\Prob(A)$ then $\Prob(B\given A)>\Prob(B)$
(assuming positive probabilities), and that $\Prob(A\given B^c)<\Prob(A)$.
\item A test for a condition has sensitivity $95\%$ and false positive rate
$2\%$. In which populations (prevalence) does a positive result give at least
a $50\%$ probability of the condition? Solve using odds.
\item Prove that $A$ is independent of itself if and only if $\Prob(A)\in\{0,1\}$.
\item Prove that if $A$ and $B$ are independent and $A$ and $C$ are independent,
it need not follow that $A$ and $B\cap C$ are independent. (Use
Example~\ref{ex:pairwise}.)
\item Show that if $A_1,\dots,A_n$ are mutually independent then
$\Prob(\bigcup A_i)=1-\prod(1-\Prob(A_i))$.
\item (Gambler's ruin warm-up) You start with $1$ unit and bet $1$ unit on fair
coin tosses until you have $0$ or $3$ units. Let $h_k$ be the probability of
reaching $3$ from $k$. Use the law of total probability on the first toss to
write equations for $h_1,h_2$ and solve them.
\item Three prisoners problem: one of three prisoners $A,B,C$ will be pardoned
at random. $A$ asks the guard to name one of the other two who will \emph{not}
be pardoned. The guard names $B$. Under the assumption that if both $B$ and $C$
are not pardoned the guard picks one at random, what is the probability that
$A$ is pardoned? Compare with the case where the guard always names $B$ when
possible.
\item (Challenge) Let $A_1,A_2,\dots$ be mutually independent with
$\Prob(A_i)=p_i$. Show that
$\Prob(\bigcap_{i=1}^\infty A_i^c)=\prod_{i=1}^\infty(1-p_i)$ using continuity
of probability, and show that this product is positive if and only if
$\sum p_i<\infty$ (assuming all $p_i<1$).
\end{enumerate}
